\chapter{Qué significa cada variable del catálogo}
\label{anexo:variables}

El Anexo~\ref{anexo:catalogo} reproduce el catálogo cerrado con su rejilla y el valor del ganador,
pero una tabla no dice qué representa cada parámetro ni por qué está ahí. Este anexo lo explica, en
el mismo orden en que el protocolo recorre las etapas, y distingue en cada caso si la variable se
optimizó, se fijó de antemano o se estresó aparte.

Conviene aclarar antes tres recuentos parecidos que conviven en el trabajo y no son intercambiables.
El sistema emplea \NumPredictores{} \emph{variables de entrada} ---los datos que los agentes miran
para puntuar una acción---. El catálogo contiene \NumParametrosCatalogo{} \emph{parámetros}, que son
los de este anexo: no describen una empresa, sino cómo se entrena y se opera el sistema. Y la
auditoría de variables reúne \NumDiagnosticos{} filas porque cruza esas entradas con todas las
series a las que el sistema calcula diagnósticos, conservando las de ambos lados.

\section{Temporal: cuándo se mira y cuánto futuro se predice}

Seis parámetros fijan el reloj del sistema. Tres de ellos se optimizaron.

\begin{description}[leftmargin=0cm,style=nextline,itemsep=0.5em]
  \item[Cadencia de observación] Cada cuántos meses se vuelve a mirar el universo, puntuar todas las
  acciones y decidir la cartera. Mirar más a menudo produce más cohortes con las que decidir, pero
  también más oportunidades de operar sobre una señal que aún no ha madurado. El ganador observa
  cada mes.
  \item[Horizonte de la etiqueta] Cuántos meses de retorno futuro aprende a ordenar el modelo. Es la
  definición de «bueno» que el sistema persigue: con horizonte de doce meses, la mejor acción es la
  que más rinde en el año siguiente, no la que sube mañana. \emph{Se optimizó}; ganó el horizonte
  anual frente al semestral.
  \item[Años de historia por ajuste] Cuánto pasado entra en cada reentrenamiento. Más historia da
  más filas pero mezcla regímenes más antiguos, quizá ya irrelevantes. \emph{Se optimizó}; ganaron
  ocho años.
  \item[Retardo de observación] Días entre el fin de un mes y el día en que el sistema mira el
  mercado. No es lo que impide usar información futura ---de eso se encarga el filtro por fecha de
  publicación, explicado en el Capítulo~\ref{chap:datos}--- sino lo que fija cuándo se mira, y con
  ello cuántos informes del trimestre ya se han publicado. \emph{Se optimizó} entre 45 y 60 días,
  con la particularidad de que el desempate favorece deliberadamente al retardo mayor, por ser el
  más conservador. Ganaron 60.
  \item[Ponderación por recencia] Si las observaciones recientes pesan más que las antiguas dentro
  del entrenamiento. \emph{Se optimizó}; ganó no ponderar, es decir, tratar por igual toda la
  ventana.
  \item[Objetivo de aprendizaje] Si el modelo aprende regresando el rango continuo o resolviendo un
  problema de ranking directo. Se fijó en la regresión del rango.
\end{description}

\section{Representación: qué ve el modelo}

Siete parámetros gobiernan qué información llega a los agentes y en qué forma. Dos se optimizaron.

\begin{description}[leftmargin=0cm,style=nextline,itemsep=0.5em]
  \item[Conjunto de variables] Cuántos bloques temáticos se activan: un subconjunto reducido de
  cinco o los once disponibles. \emph{Se optimizó}; ganó el conjunto amplio.
  \item[Momentum fundamental] Si además del valor de un ratio se añade su \emph{cambio} desde el
  informe anterior, calculado siempre con datos ya publicados. \emph{Se optimizó}; ganó incluirlo.
  \item[Régimen de mercado] Si cada observación incorpora contexto sobre el estado general del
  mercado en esa fecha, común a todas las acciones. Ganó no incluirlo, y por no inferioridad: el
  candidato no demostró ser mejor, sólo que no era peor.
  \item[Neutralización sectorial] Si cada empresa se compara con las de su propio sector antes de
  aprender, en lugar de con todo el universo. Se fijó en no neutralizar; hacerlo habría costado
  Rank-IC de forma apreciable.
  \item[Winsorización] Si los valores extremos de cada variable se recortan a un percentil antes de
  entrenar, para que un dato anómalo no domine el ajuste. Se fijó en no recortar.
  \item[Máximo de variables por agente] Cuántas variables conserva cada especialista tras la poda.
  Limitarlo reduce varianza y facilita la interpretación, a costa de perder información. El ganador
  conserva doce.
  \item[Regla de ponderación de variables] Si la importancia de cada variable la decide el propio
  modelo o una poda basada en su estabilidad fuera de muestra. Ganó la poda por estabilidad, aunque
  por muy poco: la alternativa medía ligeramente mejor pero no superaba la puerta de dominancia.
\end{description}

\section{Modelo: cómo aprende cada agente}

Cinco parámetros describen el estimador. Ninguno se optimizó en la pasada de referencia, y el
capítulo de protocolo explica por qué eso importa: el sistema no debe su capacidad predictiva al
ajuste fino del modelo.

\begin{description}[leftmargin=0cm,style=nextline,itemsep=0.5em]
  \item[Familia de modelo] Qué algoritmo usan los cinco agentes. El catálogo ofrece árboles
  potenciados o una alternativa lineal regularizada; el ganador usa árboles.
  \item[Profundidad máxima] Cuántas divisiones encadenadas puede tener cada árbol, es decir, cuán
  específicas pueden ser las interacciones que aprende. El ganador usa tres, un valor bajo.
  \item[Número de árboles] Cuántos árboles se encadenan, cada uno corrigiendo los errores del
  anterior. El ganador usa cien.
  \item[Tasa de aprendizaje] Cuánto corrige cada árbol nuevo. Valores bajos aprenden más despacio
  pero suelen generalizar mejor.
  \item[Mínimo de observaciones por hoja] Cuántas acciones debe haber al menos en cada rama final.
  Subirlo impide que el modelo construya reglas sobre un puñado de casos.
\end{description}

\section{Meta-agente: cómo se combinan los cinco}

Tres parámetros gobiernan la capa de combinación, y el primero es la decisión más consecuente de
todo el catálogo.

\begin{description}[leftmargin=0cm,style=nextline,itemsep=0.5em]
  \item[Método de combinación] Tres opciones: repartir por igual entre los cinco agentes, o aprender
  los pesos con una regresión regularizada en su variante \emph{acotada} ---donde ningún agente
  puede pesar menos del 10\,\% ni más del 50\,\%--- o \emph{libre}, sin tope alguno. \emph{Se
  optimizó}; ganó la libre, y esa elección explica que el sistema acabe concentrando más del 95\,\%
  del peso en un solo especialista.
  \item[Cohortes de historia] Cuántas cohortes trimestrales ya cerradas mira el combinador para
  estimar los pesos. El ganador usa dieciséis.
  \item[Ponderación por recencia del combinador] Si las cohortes recientes pesan más que las
  antiguas al aprender esos pesos. Se fijó en no ponderar.
\end{description}

\section{Cartera: cómo se convierte el orden en posiciones}

Doce parámetros gobiernan la gestión. Ninguno puede modificar la señal, porque se aplican después de
congelarla, y se dividen en dos grupos que conviene no mezclar.

\subsection*{Las seis que el Portfolio Study optimiza}

Son las decisiones de gestión propiamente dichas: lo que un gestor elegiría.

\begin{description}[leftmargin=0cm,style=nextline,itemsep=0.5em]
  \item[Número de posiciones] Cuántas acciones se mantienen a la vez. Pocas concentran la señal y
  hacen el resultado volátil; muchas lo estabilizan a costa de parecerse al índice. El ganador
  mantiene ocho.
  \item[Tope de efectivo] Cuánto patrimonio puede quedar sin invertir cuando ninguna candidata
  supera el umbral. Con cero, la cartera está siempre invertida al 100\,\%. El ganador tolera hasta
  un 10\,\%, aunque en la práctica nunca lo usa.
  \item[Reparto de pesos] Si todas las posiciones pesan lo mismo o el peso sigue al alfa esperado,
  con un tope de dos a uno entre la mayor y la menor. El ganador reparte en proporción al alfa.
  \item[Permanencia mínima] Cuánto debe conservarse una posición antes de poder venderla por
  cualquier motivo, expresado como fracción del horizonte del modelo. El ganador exige un horizonte
  completo, es decir, doce meses.
  \item[Suelo de cobertura] Percentil por debajo del cual una posición deja de pertenecer a la
  cartera y se vende entera. El ganador desactiva la regla.
  \item[Tolerancia a la deriva] Cuánto pueden desviarse los pesos de su objetivo antes de emitir una
  orden que los corrija. No cambia qué se tiene, sólo cada cuánto se ajusta. El ganador tolera hasta
  un 40\,\%, el valor más laxo disponible.
\end{description}

\subsection*{Las seis que se fijan y se estresan aparte}

No describen una decisión del gestor sino el entorno en el que opera, y por eso quedan fuera de la
optimización. Optimizar un coste equivaldría a elegir el mundo en el que la estrategia luce mejor.

\begin{description}[leftmargin=0cm,style=nextline,itemsep=0.5em]
  \item[Comisión] Coste fijo aplicado sobre el nocional de cada operación. Se fija en 10 puntos
  básicos y su efecto se mide en el análisis de sensibilidad a costes.
  \item[\textit{Slippage}] Diferencia adversa estimada entre el precio de referencia y el de
  ejecución, también sobre el nocional. Se fija en 20 puntos básicos.
  \item[Alfa mínimo de salida] Alfa esperado anual por debajo del cual una posición pasa a ser
  candidata a venderse.
  \item[Margen de rotación] Ventaja adicional que una candidata externa debe ofrecer por encima del
  coste de ida y vuelta para justificar una sustitución.
  \item[Endurecimiento sin fundamentales nuevos] Multiplicador que sube los umbrales en las
  observaciones en las que no ha aparecido ningún informe nuevo, cuando la señal se apoya sólo en
  precios.
  \item[Sólo venta sin fundamentales nuevos] Si en esas mismas observaciones se permite deshacer una
  posición mala pero se prohíbe comprar su reemplazo.
\end{description}

Los cuatro últimos gobiernan la mecánica de rotación, cuyo efecto ya está mediado por las seis
variables que sí se recorren. Fijarlos es una decisión declarada de antemano, y acota lo que puede
afirmarse: la rejilla responde a qué cartera es mejor \emph{dado} este entorno, no a cuál sería la
mejor sin restricciones.
